\subsection{Nagata rings}

DEFINITION 2. --- A ring A is said to be a Nagata ring if it is Noetherian and if, for every prime ideal p of A, the Noetherian integral domain A/p is Japanese (no. 1, def. 1).

Examples. --- 1) Every finitely generated algebra over a field is a Nagata ring (no. 1, example).

\begin{enumerate}
\def\labelenumi{\arabic{enumi})}
\setcounter{enumi}{1}
\item
  Every complete local Noetherian ring is a Nagata ring (no. 2, th. 2).
\item
  The ring Z is a Nagata ring (no. 1, example and corollary of prop. 1).
\item
  One can show (exerc. 30) that every finitely generated algebra over a Nagata ring is a Nagata ring.
\end{enumerate}

PROPOSITION 4. --- Let A be a Nagata ring.

\begin{enumerate}
\def\labelenumi{\alph{enumi})}
\item
  Every finite A-algebra is a Nagata ring.
\item
  For every multiplicative subset S of A, the ring \(\mathbf{S}^{-1}\mathbf{A}\) is a Nagata ring.
\item
  Let B be a finite A-algebra, \(\rho: A \to B\) the canonical homomorphism. For every prime ideal p of B, the ring B/p, which is a finite algebra over the Japanese ring \(A/\rho^{-1}(p)\) , is Japanese (no. 1, prop. 2).
\item
  Let q be a prime ideal of \(S^{-1}A\) ; then there exists a prime ideal p of A such that \(q = S^{-1}p\) . The ring ( \(S^{-1}A\) )/q is a ring of fractions of the Japanese ring A/p, hence is Japanese (no. 1, remark 2).
\end{enumerate}

THEOREM 3 (Zariski-Nagata). --- Let A be a semi-local Noetherian ring. The following conditions are equivalent:

\begin{enumerate}
\def\labelenumi{(\roman{enumi})}
\item
  A is a Nagata ring;
\item
  for every prime ideal \(\mathfrak{p}\) of A, the \(\kappa(\mathfrak{p})\) -algebra \(\kappa(\mathfrak{p}) \otimes_{\mathbf{A}} \hat{\mathbf{A}}\) is separable;
\item
  for every reduced A-algebra R, the ring \(R \otimes_{A} \hat{A}\) is reduced.
\end{enumerate}

Let us first prove the equivalence of conditions (ii) and (iii). The implication (iii) \(\Rightarrow\) (ii) is trivial; suppose conversely that A satisfies condition (ii). Then, for every A-algebra K which is a field, the ring K \(\otimes_{\mathbf{A}}\hat{\mathbf{A}}\) is reduced. Now let C be a reduced finitely generated A-algebra; the ring C, being Noetherian, is isomorphic to a subring of a finite product K \(_1\) \(\times\) \(\cdots\) \(\times\) K \(_n\) of fields (IV, § 2, no. 5, prop. 10); since \(\hat{\mathbf{A}}\) is flat over A, the ring C \(\otimes_{\mathbf{A}}\hat{\mathbf{A}}\) is isomorphic to a subring of the reduced ring \(\prod_{i}(\mathrm{K}_{i}\otimes_{\mathbf{A}}\hat{\mathbf{A}})\) , hence is reduced. Finally, let R be an arbitrary reduced A-algebra; then R is the union of the filtered family \((\mathrm{C}_{\alpha})\) of its finitely generated subalgebras, and \(R \otimes_{A} \hat{A}\) is the inductive limit of the filtered family \((\mathrm{C}_{\alpha} \otimes_{A} \hat{A})\) of reduced rings, hence is reduced.

Let us show that (ii) implies (i). Let p be a prime ideal of A; the field of fractions K of the ring A/p is identified with \(\kappa(p)\) , and the K-algebra \(K \otimes_{A/p} (\widehat{A/p})\) is identified with \(\kappa(p) \otimes_{A/p} \widehat{A}/p\widehat{A}\) , hence with \(\kappa(p) \otimes_{A} \widehat{A}\) . If \(\kappa(p) \otimes_{A} \widehat{A}\) is a \(\kappa(p)\) -separable algebra, the ring A/p is Japanese (no. 2, prop. 3).

Let us prove the implication (i) \(\Rightarrow\) (ii) by induction on dim(A). It is evident if dim(A) = 0, since then A is Artinian, hence complete. Let \(n\) be an integer \(>0\) ; consider the following hypothesis:

\((\mathbf{R}_n)\left\{
\begin{array}{ll}\text{for every local Noetherian Nagata ring C of dimension } < n\text{ and every prime}\\ \text{ideal } \mathfrak{r}\text{ of C, the ring }\kappa(\mathfrak{r})\otimes_{\mathrm{C}}\hat{\mathrm{C}}\text{ is reduced.} \end{array}
\right.\)

Let A be a semi-local Noetherian Nagata ring of dimension n, let p be a prime ideal of A, and let L be a finite extension of the field \(\kappa(\mathfrak{p})\); it suffices to prove, under hypothesis (R \(_{n}\) ), that the ring \(L \otimes_{A} \hat{A}\) is reduced. Let B denote the integral closure of A/p in L; since A/p is Japanese, B is a finite A-algebra, hence a semi-local Nagata ring (prop. 4). Let m \(_{1}\) , \ldots, m \(_{r}\) be the maximal ideals of B; the ring \(L \otimes_{A} \hat{A}\) is identified with a ring of fractions of \(B \otimes_{A} \hat{A}\), and the latter is identified with the product of the completions of the local rings B \(_{m_{i}}\) (no. 3, lemma 1). It therefore suffices to prove that, for every maximal ideal m of B, the ring B \(_{m}\) is reduced (II, § 2, no. 6, prop. 17). The ring B \(_{m}\) is local, integrally closed, and Nagata (prop. 4), and one has dim(B \(_{m}\) ) ≤ dim(B) ≤ dim(A) = n (VIII, § 1, no. 3, prop. 6 and § 2, no. 3, th. 1). Changing notations, one is reduced to proving, under hypothesis (R \(_{n}\) ), that for every local Noetherian integrally closed Nagata ring A of dimension ≤ n, the ring \(\hat{A}\) is reduced, that is to say (no. 3, lemma 2) that \(\hat{A}_{p'}\) is reduced for every prime ideal p′ ∈ Ass(\(\hat{A}\)). Since this is immediate if dim(A) = 0, we may assume dim(A) > 0. Let then x be a nonzero element of m \(_{A}\) and q' a prime ideal of \(\hat{A}\), minimal among those containing \(x\hat{A} + p'\); since \(\hat{A}_{p'}\) is identified with a ring of fractions of the ring \(\hat{A}_{q'}\), it suffices to prove that the latter is reduced (II, § 2, no. 6, prop. 17). By lemma 3, the ideal q' is associated to the \(\hat{A}\)-module \(\hat{A}/x\hat{A}\); since \(\hat{A}\) is flat over A, the inverse image q of q' in A is associated to the A-module A/xA (IV, § 2, no. 6, cor. 1 to th. 2). Since the ring A is assumed to be integrally closed, this implies that q has height 1 (VII, § 1, no. 6, prop. 10), hence that the ring \(A_{q}\) is a discrete valuation ring (loc. cit., no. 3, corollary to th. 2 and no. 6, th. 4). Since A/q is a Nagata ring of dimension < n, the ring \(\kappa(q) \otimes_{A/q} \hat{A}/q\hat{A}\) is reduced by hypothesis (R \(_{n}\) ). The ring \(\kappa(q) \otimes_{A} \hat{A}\), which is isomorphic to it, is reduced, and consequently so is the ring \(\kappa(q) \otimes_{A_{q}} \hat{A}_{q'}\), which is a ring of fractions of it. One may therefore apply Lemma 4 of Section 3 to the canonical homomorphism \(A_{q} \to \hat{A}_{q'}\), and we conclude that the ring \(\hat{A}_{q'}\) is reduced, which is what we wanted to prove. Theorem 3 is thus proved.

COROLLARY 1. --- The completion of a local and reduced Nagata ring is reduced.

It suffices indeed to set R = A in Thm. 3, (iii).

COROLLARY 2 (Chevalley). --- Let A be a reduced algebra of finite type over a field, and let p be a prime ideal of A. The completion of the local ring A \(_{p}\) is reduced.

Since A is reduced, the local ring \(A_{p}\) is reduced; moreover, A is a Nagata ring (Example 1), hence \(A_{p}\) is a Nagata ring (prop. 4), and cor. 2 follows from cor. 1, applied to the ring \(A_{p}\) .

COROLLARY 3. --- Let k be a field of characteristic 0, and let A be a local Noetherian k-algebra. For A to be a Nagata ring, it is necessary and sufficient that, for every prime ideal p of A, the ring (A/p) be reduced.

Indeed, since the fields \(\kappa(p)\) are of characteristic 0, it is equivalent to say that the algebras \(\kappa(p) \otimes_{A} \hat{A} = \kappa(p) \otimes_{A/p} (\widehat{A/p})\) are reduced or that they are separable (A, V, p.~117, th. 1), which shows that the stated condition is sufficient (th. 3, (ii) \(\Rightarrow\) (i)); it is moreover necessary (th. 3, (i) \(\Rightarrow\) (iii) with \(R = A/p\) ).

\section*{Appendix}
\phantomsection
\addcontentsline{toc}{section}{Appendix}

\refstepcounter{appendixitem}
\subsection*{\theappendixitem.\ Inductive limit of local rings}
\phantomsection
\addcontentsline{toc}{subsection}{\theappendixitem.\ Inductive limit of local rings}

Let I be a nonempty right-filtering preordered set, and let \((\mathbf{A}_{\alpha}, \varphi_{\beta\alpha})\) be an inductive system of rings relative to I. Suppose that, for every \(\alpha \in I\) , the ring \(A_{\alpha}\) is local, with maximal ideal \(m_{\alpha}\) , that the homomorphisms \(\varphi_{\beta\alpha}\) are local and flat, and that one has \(\varphi_{\beta\alpha}(m_{\alpha}) A_{\beta} = m_{\beta}\) for \(\beta \geqslant \alpha\) . Denote by A the inductive limit of the \(A_{\alpha}\) , and for every \(\alpha \in I\) , let \(\varphi_{\alpha}: A_{\alpha} \to A\) be the canonical homomorphism.

PROPOSITION 1. --- a) The ring A is local, with maximal ideal \(\mathfrak{m} = \varinjlim \mathfrak{m}_{\alpha}\) . For every \(\alpha \in \mathbf{I}\) , the homomorphism \(\varphi_{\alpha}\) is local and flat, and one has \(\varphi_{\alpha}(\mathfrak{m}_{\alpha})\) A = \(\mathfrak{m}\) .

\begin{enumerate}
\def\labelenumi{\alph{enumi})}
\setcounter{enumi}{1}
\item
  If \(\mathbf{A}_{\alpha}\) is Noetherian for every \(\alpha \in \mathbf{I}\) , then \(\mathbf{A}\) is Noetherian.
\item
  Set \(m = \varinjlim m_{\alpha}\) ; it is an ideal of A. The quotient ring A/m is the inductive limit of the fields \(A_{\alpha}/m_{\alpha}\) , it is therefore a field (A, I, p.~116, prop. 3). Moreover, every element of A - m is invertible in A: indeed, let \(x \in A - m\) ; there exists \(\alpha \in I\) and \(\xi \in A_{\alpha}\) such that \(x = \varphi_{\alpha}(\xi)\) ; one has \(\xi \notin m_{\alpha}\) , hence \(\xi\) is invertible in \(A_{\alpha}\) and x is invertible in A. Consequently, A is a local ring with maximal ideal m. Let \(\alpha \in I\) . From the relations \(\varphi_{\beta\alpha}(m_{\alpha}) A_{\beta} = m_{\beta}\) for \(\beta \geqslant \alpha\) , we deduce, by passing to the inductive limit, \(\varphi_{\alpha}(m_{\alpha}) A = m\) ; finally, the homomorphism \(\varphi_{\alpha}\) is flat by I, § 2, n° 7, prop. 9.
\item
  Let \(\hat{\mathbf{A}}\) be the separated completion of \(\mathbf{A}\) for the m-adic topology and \(\pi\) the canonical map from \(\mathbf{A}\) to \(\hat{\mathbf{A}}\) . Suppose the rings \(\mathbf{A}_{\alpha}\) are Noetherian. Fix \(\alpha \in \mathbf{I}\) and prove that the ring \(\hat{\mathbf{A}}\) is Noetherian and flat over \(\mathbf{A}_{\alpha}\) . By hypothesis, \(m_{\alpha}\) is a finitely generated ideal of \(\mathbf{A}_{\alpha}\) , hence \(m = \varphi_{\alpha}(m_{\alpha})\) . \(m\) is a finitely generated ideal of \(\mathbf{A}\) . It follows that the maximal ideal \(\hat{m}\) of \(\hat{\mathbf{A}}\) is equal to \(m\hat{\mathbf{A}}\) (III, § 2, n° 12, cor. 2 to prop. 16 and n° 13, prop. 19), hence is finitely generated. Consequently, the ring \(\hat{\mathbf{A}}\) is Noetherian (loc. cit., n° 10, cor. 5 to th. 2). On the other hand, for every \(n \in \mathbf{N}\) , the quotient \(\hat{\mathbf{A}}/\hat{m}^{n}\) is isomorphic to \(\mathbf{A}/m^{n}\) (loc. cit., n° 12, cor. 2 to prop. 16 and formula (21)), which means that \(\hat{\mathbf{A}}/\pi \circ \varphi_{\alpha}(m_{\alpha}^{n})\) \(\hat{\mathbf{A}}\) is isomorphic to \(\mathbf{A} \otimes_{\mathbf{A}_{\alpha}} (\mathbf{A}_{\alpha}/m_{\alpha}^{n})\) ; since \(\mathbf{A}\) is a flat \(\mathbf{A}_{\alpha}\)-module, the \((\mathbf{A}_{\alpha}/m_{\alpha}^{n})\) -module \(\hat{\mathbf{A}}/\pi \circ \varphi_{\alpha}(m_{\alpha}^{n})\) \(\hat{\mathbf{A}}\) is flat for every \(n \in \mathbf{N}\) . By III, § 5, n° 4, prop. 2, the \(\mathbf{A}_{\alpha}\) -module \(\hat{\mathbf{A}}\) is ideally separated for \(m_{\alpha}\) ; by loc. cit., n° 2, th. 1, the \(\mathbf{A}_{\alpha}\) -module \(\hat{\mathbf{A}}\) is therefore flat. It follows by passing to the inductive limit that \(\hat{\mathbf{A}}\) is (faithfully) flat over \(\mathbf{A}\) (I, § 2, n° 7, prop. 9), hence that \(\mathbf{A}\) is Noetherian (I, § 3, n° 5, corollary to prop. 8).
\end{enumerate}

\refstepcounter{appendixitem}
\subsection*{\theappendixitem.\ Inflation of a local ring}
\phantomsection
\addcontentsline{toc}{subsection}{\theappendixitem.\ Inflation of a local ring}

Let A be a local ring.

We denote by A{]}X{[} the local ring of the polynomial ring A{[}X{]} in the prime ideal \(m_{A}A[X]\) . This is a local ring with maximal ideal \(m_{A}A{]}X{[}\); the canonical homomorphism A → A{]}X{[} is local and flat, and the residue field of A{]}X{[} is the pure extension of \(\kappa_{A}\) generated by the class of X.

Lemma 1. — Let \(\mathbf{P} \in \mathbf{A}[\mathbf{X}]\) be a monic polynomial whose image \(\overline{\mathbf{P}}\) in \(\kappa_{\mathbf{A}}[\mathbf{X}]\) is irreducible. Then the A-algebra \(\mathbf{B} = \mathbf{A}[\mathbf{X}] / (\mathbf{P})\) is local and finite over \(\mathbf{A}\), with maximal ideal \(\mathfrak{m}_{\mathbf{A}}\mathbf{B}\); the canonical homomorphism \(\rho: \mathbf{A} \to \mathbf{B}\) is local and flat, the residue field extension \(\kappa_{\mathbf{A}} \to \kappa_{\mathbf{B}}\) is algebraic and generated by the class \(x\) of \(\mathbf{X}\) , and the minimal polynomial of \(x\) over \(\kappa_{\mathbf{A}}\) is \(\overline{\mathbf{P}}\) .

Since the polynomial P is monic, the A-module B is free of finite type (A, IV, p.~10). The ring B/m \(_{A}\) B is identified with \(\kappa_{A}[X]/(\overline{P})\) , hence is a field; the ideal m \(_{A}\) B is therefore maximal. Let q be a maximal ideal of B; then the ideal \(\rho^{-1}(q)\) is maximal (V, § 2, n° 1, prop. 1); we therefore have \(\rho^{-1}(q)=m_{A}\) whence q ⊃ m \(_{A}\) and finally q = m \(_{A}\) B. Thus the ring B is local. Lemma 1 follows immediately.

DEFINITION 1. --- Let A be a local ring. We say that an A-algebra B is an elementary swelling of A if B is isomorphic to the A-algebra A{]}X{[}, or if there exists a monic polynomial P in A{[}X{]} with irreducible image in \(\kappa_{\mathrm{A}}[\mathrm{X}]\) , such that B is isomorphic to the A-algebra A{[}X{]}/(P).

Let B be an elementary swelling of A. From the preceding, the following properties result:

\begin{enumerate}
\def\labelenumi{\alph{enumi})}
\item
  The ring B is local, and the canonical homomorphism from A to B is local and flat, and in particular injective (I, § 3, no. 5, prop. 8).
\item
  The residue field \(\kappa_{B}\) of B is a monogenic extension of the residue field \(\kappa_{A}\) of A. If \(\kappa_{B}\) is of finite degree d over \(\kappa_{A}\) , then B is a free A-module of rank d.
\item
  We have \(m_{B} = m_{A}B\) . In particular, if A is a field, then so is B. A field extension is an elementary swelling if and only if it is monogenic.
\item
  If A is Noetherian, then so is B.
\end{enumerate}

DEFINITION 2. --- Let A be a local ring. We say that an A-algebra B is a swelling of A if there exists a well-ordered set Λ having a greatest element ω, and an increasing family (B \(_{λ}\) ) \(_{λ∈Λ}\) of subalgebras of B satisfying the following conditions:

\begin{enumerate}
\def\labelenumi{\alph{enumi})}
\item
  We have \(\mathbf{B}_{\omega} = \mathbf{B}\) and the ring \(\mathbf{B}_{\lambda}\) is local for all \(\lambda \in \Lambda\) .
\item
  If \(\alpha\) is the smallest element of \(\Lambda\) , the A-algebra \(\mathbf{B}_{\alpha}\) is isomorphic to \(\mathbf{A}\) .
\item
  Let \(v \neq \alpha\) in \(\Lambda\) and let \(S_{v}\) be the set of \(\lambda \in \Lambda\) such that \(\lambda < v\) . If \(S_{v}\) has no greatest element, we have \(B_{v} = \bigcup_{\lambda \in S_{v}} B_{\lambda}\) ; if \(S_{v}\) has a greatest element \(\mu\) , then \(B_{v}\) is an elementary swelling of \(B_{\mu}\) .
\end{enumerate}

Let B be a ring and \(\rho: A \to B\) a ring homomorphism. We say that \(\rho\) is a swelling (resp. an elementary swelling) if the A-algebra defined by \(\rho\) has this property. If this is so, \(\rho\) is injective.

Examples. --- 1) Every field extension is a swelling. Indeed, let K be an extension of a field k. Endow K with a well-ordering for which 0 is the greatest element, and for \(\lambda \in \mathbf{K}\) , let \(\mathbf{K}_{\lambda}\) be the sub- \(k\) -extension of K generated by the elements \(\beta\) of K such that \(\beta < \lambda\) . The verification of the conditions \(a)\) , \(b)\) , \(c)\) , for \(k\) , K, and the family \((\mathbf{K}_{\lambda})_{\lambda \in \mathbf{K}}\) is immediate.

\begin{enumerate}
\def\labelenumi{\arabic{enumi})}
\setcounter{enumi}{1}
\tightlist
\item
  Let A be a local ring, and let I be a set of indices. Denote by A{]} \((X_{i})_{i\in I}\) {[} the local ring of the polynomial ring A \([(X_{i})_{i\in I}]\) in the prime ideal \(m_{A}A[(X_{i})_{i\in I}]\) . The A-algebra A{]} \((X_{i})_{i\in I}\) {[} is a swelling of A. Indeed, let us well-order the set I; let Λ be the well-ordered set obtained by adjoining to I a greatest element ω. For i ∈ I, identify A{]} \((X_{j})_{j<i}\) {[} to a subalgebra B \(_{i}\) of B = A{]} \((X_{i})_{i\in I}\) , and set B \(_{\omega}\) = B. The family (B \(_{\lambda}\) ) \(_{\lambda\in\Lambda}\) satisfies conditions a), b), c).
\end{enumerate}

Remark. --- With the notations of Def. 2, the ring \(\mathbf{B}_{\mu}\) is a swelling of \(\mathbf{B}_{\lambda}\) when \(\lambda \leqslant \mu\) .

PROPOSITION 2. --- Let A be a local ring and B a swelling of A.

\begin{enumerate}
\def\labelenumi{\alph{enumi})}
\item
  The ring B is local, and we have \(\mathfrak{m}_{\mathrm{A}}\mathbf{B} = \mathfrak{m}_{\mathrm{B}}\) .
\item
  The A-algebra B is faithfully flat.
\item
  The canonical homomorphism
\end{enumerate}

\[
\gamma_ {\mathrm{B}}: \operatorname{gr} (\mathrm{A}) \otimes_ {\kappa_ {\mathrm{A}}} \kappa_ {\mathrm{B}} \rightarrow \operatorname{gr} (\mathrm{B})
\]

is bijective.

\begin{enumerate}
\def\labelenumi{\alph{enumi})}
\setcounter{enumi}{3}
\tightlist
\item
  If A is Noetherian, then so is B, and the Hilbert-Samuel series (VIII, § 4, no. 3) of A and B are equal.
\end{enumerate}

Let \((\mathbf{B}_{\lambda})_{\lambda\in\Lambda}\) a family of subalgebras of B satisfying conditions a), b), and c) of def. 2.

Let \(\Lambda'\) be the set of indices \(\lambda \in \Lambda\) such that, for every \(\mu \leqslant \lambda\) in \(\Lambda\) , the A-algebra \(B_{\mu}\) is local and faithfully flat, and \(m_{B_{\mu}} = m_A B_{\mu}\) . Suppose \(\Lambda' \neq \Lambda\) and let v be the smallest element of \(\Lambda - \Lambda'\) . We have \(\alpha \in \Lambda'\) , whence v \(\neq \alpha\) . \(S_{v}\) is contained in \(\Lambda'\) . If \(S_{v}\) has no greatest element, we have \(B_{v} = \bigcup_{\lambda \in S_{v}} B_{\lambda}\) and v belongs to \(\Lambda'\)

  by prop. 1 of no. 1. If \(S_v\) has a greatest element \(\mu\) , one has \(\mu \in \Lambda'\) and \(B_v\) is an elementary swelling of \(B_\mu\) : we again have \(v \in \Lambda'\) by the remarks following def. 1, whence a contradiction.

When A is Noetherian, one proves in an analogous manner that the set \(\Lambda''\) of indices \(\lambda \in \Lambda\) such that the ring \(\mathbf{B}_{\lambda}\) is Noetherian is equal to \(\Lambda\) .

We therefore have \(\omega\in\Lambda'\) , whence the assertions \(a\) and \(b\) ). When A is Noetherian, we have \(\omega\in\Lambda''\) , hence B = B \(_{\omega}\) is Noetherian.

Assertion c) follows from a), b), and th. 1 of III, § 5, no. 2. Suppose A (hence B) is Noetherian; since we have

\[
\left[ \mathfrak {m} _ {\mathrm{B}} ^ {n} / \mathfrak {m} _ {\mathrm{B}} ^ {n + 1}: \kappa_ {\mathrm{B}} \right] = \left[ \mathfrak {m} _ {\mathrm{A}} ^ {n} / \mathfrak {m} _ {\mathrm{A}} ^ {n + 1}: \kappa_ {\mathrm{A}} \right]
\]

for every \(n \in \mathbf{N}\) , the Hilbert-Samuel series of A and B are equal.

COROLLARY. --- Suppose A is Noetherian.

\begin{enumerate}
\def\labelenumi{\alph{enumi})}
\item
  We have \(\dim (\mathbf{A}) = \dim (\mathbf{B})\)
\item
  Suppose A is regular, and let \((x_{1},\ldots ,x_{n})\) be a system of coordinates of A. Then B is regular and the sequence \((x_{1}1_{\mathrm{B}},\dots,x_{n}1_{\mathrm{B}})\) is a system of coordinates of B.
\end{enumerate}

This follows from prop. 1 of VIII, § 5, no. 1.

PROPOSITION 3. --- Let A, B, C be three local rings and \(u:A\to B\) , \(v:B\to C\) two swellings. Then \(v\circ u\) is a swelling.

Let \((\mathrm{B}_{\lambda})_{\lambda\in\Lambda}\) and \((\mathrm{C}_{\mu})_{\mu\in\mathbf{M}}\) be families of sub-A-algebras of B and of sub-B-algebras of C respectively, having properties a), b), c) of def. 2. On the set N, the sum of \(\Lambda\) and M, consider the order relation inducing on \(\Lambda\) and M the given orders and such that one has \(\lambda<\mu\) for \(\lambda\in\Lambda\) , \(\mu\in M\) . This is a well-ordering relation. For \(\lambda\in\Lambda\subset N\) , set \(C_{\lambda}=v(\mathrm{B}_{\lambda})\) . Then the family \((\mathrm{C}_{\nu})_{\nu\in\mathbf{N}}\) satisfies conditions a), b), c) of def. 1 relative to the A-algebra C.

THEOREM 1. --- Let \(f:A\to A'\) be a surjective local homomorphism of local rings and let B' be a swelling of A'. There exists a swelling B of A and an isomorphism of A-algebras from B \(\otimes_{A}A'\) onto B'.

\begin{enumerate}
\def\labelenumi{\Alph{enumi})}
\tightlist
\item
  Suppose that B' is an elementary swelling of A'. We distinguish two cases: 1) If B' is finite over A', choose an isomorphism of A'-algebras \(\varphi: A'[X]/(P') \to B'\) , where P' ∈ A'{[}X{]} is a monic polynomial whose image is irreducible in \(\kappa_{A'}[X]\) . Choose a monic polynomial P ∈ A{[}X{]} whose image in A'{[}X{]} is P'. It is necessarily irreducible modulo the maximal ideal of A. Then set B = A{[}X{]}/(P). The A-algebra B is an elementary swelling of A and \(\varphi\) induces an isomorphism of A-algebras from \(B \otimes_{A} A'\) onto B'.
\end{enumerate}

\begin{enumerate}
\def\labelenumi{\arabic{enumi})}
\setcounter{enumi}{1}
\tightlist
\item
  If B' is not finite over A', choose an isomorphism of A'-algebras \(\psi : A']X[\to B'\) . Put \(B = A]X[\) . The A-algebra B is an elementary swelling of A, and B \(\otimes_{A} A'\) is canonically isomorphic to A'{]}X{[}. Consequently \(\psi\) induces an isomorphism of A-algebras from B \(\otimes_{A} A'\) onto B'.
\end{enumerate}

\begin{enumerate}
\def\labelenumi{\Alph{enumi})}
\setcounter{enumi}{1}
\tightlist
\item
  Let us pass to the general case. Let \((\mathbf{B}_{\lambda}^{\prime})_{\lambda\in\Lambda}\) be a family of sub-\(A^{\prime}\)-algebras of \(B^{\prime}\) having, relative to \(A^{\prime}\) and \(B^{\prime}\) , the properties a), b), c) of def. 2. We shall define by transfinite induction an inductive system \((\widetilde{\mathbf{B}}_{\lambda}, i_{\mu\lambda})\) relative to \(\Lambda\) of local rings and injective local homomorphisms, and isomorphisms \(u_{\lambda}: \widetilde{B}_{\lambda} \otimes_{A} A^{\prime} \to B_{\lambda}^{\prime}\) such that, for \(\lambda \leqslant \mu\) , \(u_{\mu} \circ (i_{\mu\lambda} \otimes \operatorname{Id}_{A^{\prime}}) \circ u_{\lambda}^{-1}\) is the canonical injection of \(B_{\lambda}^{\prime}\) in \(B_{\mu}^{\prime}\) .
\end{enumerate}

If \(\alpha\) is the smallest element of \(\Lambda\) , we set \(\tilde{\mathbf{B}}_{\alpha} = \mathbf{A}\) , \(i_{\alpha \alpha} = \mathrm{Id}_{\mathbf{A}}\) and we take for \(u_{\alpha}\) the canonical isomorphism \(\mathbf{A} \otimes_{\mathbf{A}} \mathbf{A}' \to \mathbf{A}'\) .

Let \(v \in \Lambda\) , and suppose \(\tilde{\mathbf{B}}_{\lambda}\) , \(u_{\lambda}\) and \(i_{\mu \lambda}\) have been constructed when \(\lambda \leqslant \mu < v\) . Let \(S_v\) be the set of elements \(\varepsilon\) of \(\Lambda\) such that \(\varepsilon < v\) . If \(S_v\) has no greatest element, one takes for \(\tilde{\mathbf{B}}_v\) the inductive limit of the \(\tilde{\mathbf{B}}_{\lambda}\) for \(\lambda \in S_v\) , for \(u_v\) the composed isomorphism \(\tilde{\mathbf{B}}_v \otimes_A A' \to \varinjlim (\tilde{\mathbf{B}}_{\lambda} \otimes_A A') \to \varinjlim B'_\lambda \to B'_v\) , and for \(i_{v\lambda}\) , when \(\lambda \in S_v\) the canonical map from \(\tilde{\mathbf{B}}_{\lambda}\) to \(\tilde{\mathbf{B}}_v\) . If \(S_v\) has a greatest element \(\mu\) , then \(B'_v\) is an elementary swelling of \(B'_\mu\) . By A), there exists an elementary swelling \(i_{\nu \mu}:\tilde{\mathbf{B}}_{\mu}\to \tilde{\mathbf{B}}_{\nu}\) and an isomorphism of \(\tilde{\mathbf{B}}_{\mu}\) -algebras from \(\tilde{\mathbf{B}}_{\nu}\otimes_{\tilde{\mathbf{B}}_{\mu}}\mathbf{B}_{\mu}^{\prime}\) onto \(\mathbf{B}_{\nu}^{\prime}\) . Take for \(u_{\nu}\) the isomorphism of A-algebras composed

\[
\widetilde {\mathbf {B}} _ {\nu} \otimes_ {\mathrm{A}} \mathrm{A} ^ {\prime} \rightarrow \widetilde {\mathbf {B}} _ {\nu} \otimes_ {\widetilde {\mathbf {B}} _ {\mu}} (\widetilde {\mathbf {B}} _ {\mu} \otimes_ {\mathrm{A}} \mathrm{A} ^ {\prime}) \rightarrow \widetilde {\mathbf {B}} _ {\nu} \otimes_ {\widetilde {\mathbf {B}} _ {\mu}} \mathrm{B} _ {\mu} ^ {\prime} \rightarrow \mathrm{B} _ {\nu} ^ {\prime}
\]

and for \(i_{\nu \lambda}\) , when \(\lambda \in S_{v}\) , the homomorphism \(i_{\nu \mu} \circ i_{\mu \lambda}\) .

Then set \(\mathbf{B} = \tilde{\mathbf{B}}_{\omega}\) and, for all \(\lambda \in \Lambda\) , denote by \(\mathbf{B}_{\lambda}\) the image of \(\tilde{\mathbf{B}}_{\lambda}\) by the canonical injection \(\tilde{\mathbf{B}}_{\lambda} \to \mathbf{B}\) . The family \((\mathbf{B}_{\lambda})_{\lambda \in \Lambda}\) satisfies the conditions \(a)\) , \(b)\) , \(c)\) of def. 2, and \(\mathbf{B}\) is a swelling of A. On the other hand, the homomorphism \(u_{\omega}\) is an A'-isomorphism of \(\mathbf{B} \otimes_{A} \mathbf{A}'\) onto \(\mathbf{B}'\) .

COROLLARY. --- Let A be a local ring and K an extension of its residue field \(\kappa_{\mathrm{A}}\). There exists a local ring B and a swelling \(\mathbf{A} \to \mathbf{B}\) such that the \(\kappa_{\mathrm{A}}\) -algebra \(\kappa_{\mathrm{B}}\) is isomorphic to K.

Indeed, the homomorphism \(\kappa_{\mathbf{A}} \to \mathbf{K}\) is a swelling (example 1). Applying Thm. 1 with \(\mathbf{A}' = \kappa_{\mathbf{A}}\) and \(\mathbf{B}' = \mathbf{K}\) , we obtain the existence of a swelling B of A and an A-isomorphism of \(\mathbf{B} / \mathfrak{m}_{\mathbf{A}}\mathbf{B}\) onto K, whence the corollary.

\refstepcounter{appendixitem}
\subsection*{\theappendixitem.\ Existence of p-rings}
\phantomsection
\addcontentsline{toc}{subsection}{\theappendixitem.\ Existence of p-rings}

PROPOSITION 4. --- Let \(p\) be a prime number, \(k\) be a field of characteristic \(p\) , and let \(n\) be an integer \(\geqslant 1\) , or \(+\infty\) . There exists a \(p\) -ring (§ 2, no. 1, def. 1) of length \(n\) whose residue field is isomorphic to \(k\) .

One may consider \(k\) as an extension of the residue field \(\mathbf{Z}/p\mathbf{Z}\) of the local ring \(\mathbf{Z}_{(p)}\) . By the corollary of Th. 1, there exists a local ring \(\mathbf{B}\) , a swelling of \(\mathbf{Z}_{(p)}\), such that \(\kappa_{\mathrm{B}}\) is isomorphic to \(k\) . The local ring \(\mathbf{Z}_{(p)}\) is regular and \(\{p\}\) is a system of coordinates of \(\mathbf{Z}_{(p)}\) . By the corollary of prop. 2 of no. 2, the ring \(\mathbf{B}\) is regular and \(\{p1_{\mathrm{B}}\}\) is a system of coordinates of \(\mathbf{B}\) . In other words, \(\mathbf{B}\) is a discrete valuation ring, with maximal ideal \(p\mathbf{B}\) . The completion \(\mathbf{C}\) of \(\mathbf{B}\) is then a \(p\) -ring of infinite length and the residue field \(\kappa_{\mathrm{C}}\) is isomorphic to \(\kappa_{\mathrm{B}}\) , hence to \(k\) . Moreover, for every integer \(n \geqslant 1\) , C/ \(p^{n}\) C is a \(p\) -ring of length \(n\) , with residue field isomorphic to \(\kappa_{\mathrm{C}}\) , hence to \(k\) .

\specialchapter{Exercises}
\setcounter{exercisesection}{0}
